\chapter{Circumcision}

\section*{What Is This Surgery?}
Circumcision is a common surgical procedure involving the complete or partial removal of the foreskin, also known as the prepuce, which is the retractable fold of skin covering the glans penis. This procedure can be performed at various ages, most frequently in the neonatal period, but also in adolescence or adulthood for specific medical indications. Its practice is deeply rooted in cultural, religious, and traditional beliefs in many societies worldwide, while in others, it is primarily undertaken for perceived health benefits or therapeutic reasons. The primary anatomical goal is to permanently expose the glans penis, thereby altering its hygiene and potentially its susceptibility to certain conditions, making it a significant procedure in both elective and medically necessary contexts.

\section*{Alternative Names}
\begin{itemize}
  \item Male ritual circumcision
  \item Neonatal circumcision
  \item Adult circumcision
  \item Preputial excision
  \item Posthectomy
  \item Penile foreskin removal
\end{itemize}

\section*{Relevant Surgical Anatomy}
The relevant surgical anatomy for circumcision primarily involves the penis, specifically the prepuce (foreskin) and its relationship to the glans penis and penile shaft. The prepuce is a double-layered fold of skin that covers the glans penis. It consists of an outer skin layer continuous with the penile shaft skin and an inner mucosal layer that is continuous with the skin of the glans at the coronal sulcus. The frenulum is a small band of tissue connecting the ventral aspect of the glans to the inner surface of the foreskin, containing a small artery and nerve. Deep to the skin layers are the dartos fascia, a thin layer of smooth muscle and connective tissue, and then Buck's fascia, which encases the corpora cavernosa and corpus spongiosum.

The arterial supply to the foreskin is primarily from branches of the internal pudendal artery, specifically the dorsal penile arteries and superficial external pudendal arteries. Venous drainage occurs via the superficial dorsal vein of the penis, which drains into the saphenous system, and the deep dorsal vein, which drains into the prostatic plexus. Innervation is provided by branches of the pudendal nerve, including the dorsal nerves of the penis, which supply sensation to the glans and foreskin. Lymphatic drainage from the penis, including the foreskin, primarily flows to the superficial inguinal lymph nodes. Key surgical landmarks include the coronal sulcus, which is the groove separating the glans from the shaft, and the frenulum, which must be carefully managed during excision.

\section*{Surgical Principle \& Rationale}
The fundamental principle of circumcision is the surgical excision of the preputial skin, which normally covers the glans penis, to achieve permanent exposure of the glans. This procedure aims to eliminate the potential for conditions such as pathological phimosis (inability to retract the foreskin), paraphimosis (foreskin trapped behind the glans), and recurrent balanoposthitis (inflammation of the glans and foreskin). By removing the foreskin, the moist, warm sub-preputial environment is eliminated, which is thought to reduce bacterial and fungal colonization and improve local hygiene.

From a prophylactic standpoint, the exposed glans is believed to be less susceptible to certain infections, including urinary tract infections (UTIs) in infancy and sexually transmitted infections (STIs) such as HIV, human papillomavirus (HPV), and herpes simplex virus (HSV) in adulthood, although the extent of this protection remains a subject of ongoing research and debate. The procedure also serves as a definitive treatment for various pathological conditions of the foreskin that are refractory to conservative management, such as balanitis xerotica obliterans (BXO). The technical goals are to achieve a clean, aesthetically acceptable surgical margin with adequate hemostasis, ensuring proper healing, preserving penile function, and minimizing complications.

\section*{History \& Surgical Evolution}
Circumcision is one of the oldest surgical procedures known to humanity, with evidence of its practice dating back thousands of years. Ancient Egyptian tomb drawings from as early as 2300 BC depict the procedure, suggesting its significance in that civilization, likely for religious or ritualistic purposes. Its origins are also deeply intertwined with various religious and cultural traditions, particularly within Judaism, where it is a covenantal rite (Brit Milah) performed on the eighth day of life, and in Islam, where it is a widely practiced tradition for both religious and cultural reasons.

In the Western world, particularly in English-speaking countries, the practice gained significant medical traction in the 19th and early 20th centuries. It was promoted for a variety of perceived health benefits, including improved hygiene, prevention of venereal diseases, and even as a treatment for conditions ranging from epilepsy to masturbation. This era saw the standardization of various surgical techniques and the development of specialized instruments such as the Gomco clamp (developed by Dr. Hiram S. Gomco in the 1930s) and the Plastibell device (introduced in the 1950s). Modern techniques, such as the sleeve resection, prioritize precise cosmetic outcomes and meticulous hemostasis, building upon centuries of surgical evolution to enhance safety and efficacy.

\section*{Clinical Indications}
Circumcision is indicated for a variety of reasons, ranging from elective to medically necessary:
\begin{itemize}
  \item **Cultural or Religious Preference:** The most common reason for neonatal circumcision in many parts of the world, driven by long-standing traditions and beliefs.
  \item **Pathological Phimosis:** A condition where the foreskin is too tight to be retracted over the glans, causing pain, difficulty with hygiene, recurrent infections, or urinary obstruction. This is often diagnosed after conservative measures like topical steroids have failed.
  \item **Paraphimosis:** A urological emergency where the retracted foreskin becomes trapped behind the glans, leading to swelling, pain, and potential vascular compromise of the glans. Circumcision is typically performed after manual reduction has failed or to prevent recurrence.
  \item **Recurrent Balanoposthitis:** Chronic or recurrent inflammation of the glans and foreskin, often due to poor hygiene, fungal infections, or diabetes, unresponsive to conservative treatment with antibiotics or antifungals.
  \item **Balanitis Xerotica Obliterans (BXO):** A chronic inflammatory skin condition that can cause progressive scarring, tightening, and whitening of the foreskin and meatus, often leading to severe phimosis and requiring surgical excision.
  \item **Suspected Penile Cancer:** In cases where a lesion on the foreskin is suspicious for malignancy, circumcision may be performed for diagnostic biopsy and definitive treatment, often as part of a wider local excision.
  \item **Recurrent Urinary Tract Infections (UTIs) in Infants:** Recommended in some high-risk male infants, particularly those with vesicoureteral reflux, hydronephrosis, or other urinary tract anomalies, where the foreskin may harbor bacteria.
  \item **Reduction of Sexually Transmitted Infection (STI) Transmission:** While not a primary indication for elective circumcision in many developed countries, it is recognized as a factor in reducing the risk of HIV, HPV, and HSV in certain high-prevalence populations, as supported by randomized controlled trials.
\end{itemize}

\section*{Clinical Positioning}
Circumcision can serve as both a primary and a secondary procedure, depending on the clinical context and the underlying condition. It is considered a primary procedure when performed electively in the neonatal period for cultural or religious reasons, as it is the initial and definitive intervention. Similarly, for acute conditions like paraphimosis, it is often the primary surgical intervention after attempts at manual reduction have failed, or to prevent recurrence. In cases of pathological phimosis, circumcision is frequently the primary surgical intervention after conservative measures, such as topical steroid creams, have proven ineffective or are deemed inappropriate due to severity or patient preference.

Conversely, circumcision may be a secondary or salvage procedure when it is undertaken after other treatments have proven ineffective or when it is part of a broader management strategy. For instance, in recurrent balanoposthitis, it is typically considered after trials of antifungal creams, antibiotics, and improved hygiene have not resolved the issue. In cases of balanitis xerotica obliterans (BXO), it may be performed after initial topical steroid therapy has failed to control the progressive scarring. Similarly, in cases of suspected penile cancer, it may be a secondary procedure following initial biopsy or as part of a more extensive oncological resection. Its role is therefore dictated by the urgency, etiology, and chronicity of the underlying condition, ranging from an elective prophylactic measure to a therapeutic intervention for refractory disease.

\section*{Surgical Technique \& Procedure}
The technique for circumcision varies significantly between neonates and older males, primarily due to anatomical differences, the need for different anesthetic approaches, and the desired cosmetic outcome.

In **neonates**, the procedure is typically performed under local anesthesia, often a dorsal penile nerve block combined with a ring block, to minimize discomfort. The infant is positioned supine, and the penis is prepped and draped. One of three common devices is usually employed:
\begin{itemize}
  \item **Gomco Clamp:** After separating the foreskin from the glans (adhesiolysis), a bell-shaped device is placed over the glans, and the foreskin is drawn over the bell. The Gomco clamp is then applied, compressing the foreskin against the bell for 5-10 minutes to achieve hemostasis before the foreskin is excised with a scalpel.
  \item **Plastibell:** A plastic bell is placed over the glans, and the foreskin is tied tightly around its groove with a ligature. The excess foreskin distal to the ligature is then trimmed. The plastic bell typically remains in place for 5-7 days, falling off spontaneously as the ligated tissue necroses.
  \item **Mogen Clamp:** This device is used for a rapid excision. The foreskin is pulled through a slot in the clamp, ensuring the glans is protected. The clamp is then closed, and the foreskin is excised with a scalpel along the distal edge of the clamp.
\end{itemize}

In **older children and adults**, the procedure is usually performed under local anesthesia (dorsal penile block, ring block, or penile block with local infiltration), regional anesthesia (spinal or epidural), or general anesthesia, depending on patient preference, age, and complexity. The patient is positioned supine. After sterile preparation and draping, the common surgical techniques include the dorsal slit technique and the sleeve resection technique. The sleeve resection is generally preferred for its superior cosmetic outcome.

**Key Steps for Adult Circumcision (Sleeve Resection Technique):**
\begin{itemize}
  \item **Anesthesia:** Local infiltration with lidocaine/bupivacaine, often combined with a dorsal penile nerve block, or general anesthesia.
  \item **Marking:** The desired lines of excision are marked on the foreskin. Typically, an inner line is marked just proximal to the coronal sulcus (e.g., 0.5-1 cm) and an outer line more proximally on the penile shaft, ensuring adequate skin removal without tension.
  \item **Dorsal Slit (Optional):** A small dorsal incision may be made through the foreskin to facilitate retraction and separation of any preputial adhesions to the glans.
  \item **Dissection:** Two circumferential incisions are made along the marked lines. The intervening sleeve of foreskin is then carefully dissected from the underlying dartos fascia and penile shaft, ensuring preservation of the neurovascular bundles and frenular artery.
  \item **Hemostasis:** Bleeding vessels, particularly those in the frenulum and along the cut edges, are meticulously coagulated using electrocautery or ligated with fine absorbable sutures (e.g., 5-0 polyglactin).
  \item **Closure:** The inner (mucosal) and outer (skin) edges are approximated and sutured together using fine, absorbable sutures (e.g., 4-0 or 5-0 chromic gut or polyglactin) in an interrupted or continuous fashion, creating a neat scar just proximal to the coronal sulcus.
  \item **Dressing:** A sterile dressing, often petroleum jelly-impregnated gauze or a non-adherent dressing, is applied to the wound to provide compression and protect the surgical site.
\end{itemize}

\section*{Preoperative Workup \& Preparation}
Preparation for circumcision varies based on the patient's age, medical status, and the type of anesthesia planned. Thorough preoperative assessment is crucial to ensure patient safety and optimize outcomes.

For **neonates**:
\begin{itemize}
  \item **NPO Status:** Typically kept NPO (nothing by mouth) for 1-2 hours prior to the procedure to reduce the risk of aspiration, though this can vary by institutional policy.
  \item **Vitamin K Administration:** It is imperative to ensure the infant has received prophylactic vitamin K at birth to prevent hemorrhagic disease of the newborn, as this is crucial for normal coagulation.
  \item **Parental Consent:** A detailed discussion with parents regarding the procedure, its risks, benefits, and alternatives, followed by obtaining informed consent, is mandatory.
  \item **Pain Management:** Consideration of non-pharmacological pain relief (e.g., sucrose solution, pacifier) in addition to local anesthesia.
\end{itemize}

For **older children and adults**:
\begin{itemize}
  \item **Medical History and Physical Examination:** A thorough assessment to identify any underlying medical conditions (e.g., diabetes, cardiac disease), allergies, or bleeding disorders. Examination of the penis for any anomalies (e.g., hypospadias, chordee) is critical.
  \item **Laboratory Tests:** Routine blood tests are usually not required for healthy individuals undergoing local anesthesia. However, for general anesthesia, if there's a history of bleeding, or if anticoagulants are used, a complete blood count (CBC) and coagulation studies (PT/INR, PTT) may be ordered.
  \item **Anesthesia Consultation:** For general or regional anesthesia, an anesthesia pre-assessment is necessary to evaluate fitness for anesthesia and discuss options.
  \item **NPO Status:** If general anesthesia is planned, the patient must fast for at least 6-8 hours prior to surgery (NPO after midnight).
  \item **Medication Review:** Patients should inform their surgeon about all medications, especially anticoagulants (e.g., warfarin, aspirin, NSAIDs), which may need to be temporarily discontinued or adjusted according to a specific protocol to minimize bleeding risk.
  \item **Antibiotics:** Prophylactic antibiotics are generally not routinely indicated for elective circumcision but may be considered in specific high-risk cases, such as immunocompromised patients or those with prosthetic implants.
  \item **Informed Consent:** Comprehensive discussion of the procedure, potential risks, benefits, alternatives, and expected recovery, followed by signed informed consent.
  \item **Hygiene:** Patients are usually advised to shower and clean the genital area thoroughly on the day of surgery.
\end{itemize}

\section*{Contraindications \& Precautions}
Circumcision, while common, has specific contraindications and requires precautions to ensure patient safety and optimal outcomes. Careful patient selection is paramount.

**Absolute Contraindications:**
\begin{itemize}
  \item **Hypospadias or Epispadias:** Congenital conditions where the urethral opening is not at the tip of the penis. The foreskin is often crucial for reconstructive surgery to correct these anomalies, and its removal would complicate future repair.
  \item **Chordee:** A congenital curvature of the penis, often associated with hypospadias, where the foreskin may be needed for repair or to provide skin for straightening procedures.
  \item **Ambiguous Genitalia:** In cases where the sex of the infant is not clearly defined, circumcision should be deferred until a definitive diagnosis and treatment plan are established by a multidisciplinary team.
  \item **Severe Bleeding Disorders:** Uncontrolled coagulopathies (e.g., severe hemophilia, severe thrombocytopenia) pose a significant and unacceptable risk of hemorrhage.
  \item **Acute Local Infection:** Active infection of the glans or foreskin (e.g., severe balanitis, cellulitis, herpes simplex lesions) should be treated and resolved before elective circumcision to prevent spread of infection and impaired wound healing.
  \item **Other Significant Penile Anomalies:** Any other significant congenital anomaly of the penis where the foreskin might be required for future reconstructive surgery or where its removal could exacerbate the anomaly.
\end{itemize}

**Relative Contraindications and Precautions:**
\begin{itemize}
  \item **Prematurity or Unstable Neonate:** Circumcision should be delayed until the infant is medically stable, has reached an appropriate weight (typically >2500g), and is discharged from the hospital, as premature infants have higher risks of complications.
  \item **Significant Systemic Illness:** Any acute or chronic illness that would increase surgical or anesthetic risk (e.g., severe cardiac disease, uncontrolled diabetes) should be optimized before elective surgery.
  \item **Minor Coagulopathy:** Mild bleeding disorders or patients on antiplatelet/anticoagulant medications may require careful preoperative management, such as temporary cessation of medication or factor replacement, and meticulous intraoperative hemostasis.
  \item **Parental Indecision or Lack of Informed Consent:** The procedure should not proceed without clear, informed, and voluntary consent from the legal guardians or the patient themselves.
  \item **Undiagnosed Penile Lesions:** If a suspicious lesion is present on the foreskin or glans, it should be thoroughly evaluated and biopsied before proceeding with circumcision, as the foreskin may be needed for diagnostic purposes or reconstruction.
  \item **Anatomical Variations:** Careful assessment for a short frenulum, buried penis, or other minor variations that might influence the surgical technique or require specific modifications to prevent complications.
\end{itemize}

\section*{Complications \& Risks}
While generally considered a safe procedure, circumcision carries potential risks and complications, which can be categorized as general surgical risks and procedure-specific complications. The overall complication rate is low, typically ranging from 0.2\% to 3\%.

**General Surgical Risks:**
\begin{itemize}
  \item **Bleeding/Hemorrhage:** The most common complication, ranging from minor oozing requiring pressure to significant blood loss requiring re-exploration or transfusion. This is more common in neonates and patients with coagulopathies.
  \item **Infection:** Localized wound infection (cellulitis, abscess) or, rarely, systemic infection (sepsis). Risk factors include poor hygiene, inadequate sterile technique, and immunocompromised status.
  \item **Anesthesia Complications:** Risks associated with local, regional, or general anesthesia, including allergic reactions, respiratory depression, cardiovascular events, or nerve injury from injection.
  \item **Pain:** Postoperative pain, which is usually manageable with analgesics but can be severe if not adequately addressed.
  \item **Scarring:** Formation of hypertrophic or keloid scars, though rare on the penis.
\end{itemize}

**Procedure-Specific Risks:**
\begin{itemize}
  \item **Removal of Too Much or Too Little Skin:** Can lead to an undesirable cosmetic outcome, discomfort, or recurrence of symptoms (if too little skin is removed, e.g., persistent phimosis). Too much skin removal can cause tension and pain, especially during erection.
  \item **Meatal Stenosis:** Narrowing of the urethral opening, often due to irritation, inflammation, or scarring of the glans post-circumcision, which can cause difficulty with urination and may require further surgical correction (meatotomy).
  \item **Glans Injury:** Accidental trauma to the glans penis, ranging from superficial abrasions or lacerations to partial amputation (extremely rare but severe, particularly with clamp techniques if not properly applied).
  \item **Poor Cosmetic Outcome:** Unsatisfactory appearance of the penis, including uneven skin edges, excessive scarring, asymmetry, or an irregular scar line.
  \item **Skin Bridge Formation:** Adhesion of the penile shaft skin to the glans, creating a bridge of skin that can trap debris, cause discomfort, or interfere with hygiene, often requiring minor surgical revision.
  \item **Inclusion Cysts:** Small, benign cysts that can form along the suture line due to trapped epithelial cells during wound closure.
  \item **Urinary Retention:** Temporary difficulty or inability to urinate post-procedure, more common in adults due to pain or swelling.
  \item **Frenular Injury or Short Frenulum:** Damage to the frenulum during excision can lead to persistent bleeding or pain. Inadequate excision or a naturally short frenulum can cause pain or curvature during erection.
  \item **Nerve Damage:** Rare, but potential for injury to superficial sensory nerves, leading to altered sensation (hypoesthesia or hyperesthesia) of the glans or penile shaft.
\end{itemize}

\section*{Postoperative Recovery Timeline}
The postoperative recovery timeline for circumcision varies depending on the patient's age and the technique used, but generally involves a period of wound healing and gradual return to normal activities.

Immediately after circumcision, the patient is monitored in a recovery area (PACU for adults, or with parents for neonates) to ensure stable vital signs, adequate pain control, and absence of immediate complications like significant bleeding. For neonates, a petroleum jelly-impregnated gauze dressing is typically applied to the wound site, which should be changed with each diaper change for the first 3-5 days. Parents are instructed to keep the area clean and dry, and to apply petroleum jelly directly to the glans to prevent it from sticking to the diaper. Mild swelling, redness, and a yellowish exudate (fibrin) are normal parts of the healing process in the first week. Complete healing of the neonatal wound typically occurs within 1-2 weeks.

For older children and adults, a light dressing is usually applied, and oral analgesics (e.g., paracetamol, ibuprofen, or mild opioids) are prescribed for pain management. Patients are advised to avoid strenuous activity, heavy lifting, and sexual activity for 4-6 weeks to allow for complete wound healing and prevent dehiscence or bleeding. Swelling and bruising of the penis are common and typically resolve within 1-2 weeks. Absorbable sutures usually dissolve on their own within 2-4 weeks, though some may take longer. Complete epithelialization of the wound typically takes 2-4 weeks, but the final cosmetic appearance and scar maturation may take several months to fully develop. Patients are encouraged to maintain good hygiene, gently clean the area, and wear loose-fitting clothing to minimize irritation.

\section*{Surgical Findings \& Pathology}
During circumcision, the surgeon meticulously documents the intraoperative findings, which are crucial for understanding the underlying pathology and guiding postoperative care. These findings may include the degree of phimosis (e.g., pinhole meatus, non-retractile foreskin), presence and severity of preputial adhesions to the glans, evidence of balanitis (inflammation, erythema, discharge), or the characteristic white, sclerotic plaques of balanitis xerotica obliterans (BXO). In cases of paraphimosis, the surgeon notes the extent of glans edema and vascular compromise. For suspected malignancy, the location and appearance of any suspicious lesions are recorded.

The resected foreskin is typically sent for pathological examination, especially in adults or when there is a medical indication beyond elective neonatal circumcision. The pathology report provides definitive diagnostic information. For phimosis, it may show chronic inflammation and fibrosis of the preputial tissue. In cases of recurrent balanoposthitis, the report often confirms chronic inflammatory changes. For BXO, the characteristic histopathological features include epidermal atrophy, hyperkeratosis, and a dense band of subepithelial sclerosis. If a suspicious lesion was excised, the pathology report will confirm the presence or absence of malignancy (e.g., squamous cell carcinoma, carcinoma in situ) and assess the surgical margins, which is critical for oncological management. This pathological confirmation is vital for guiding further treatment and prognosis.

\section*{Expected Postoperative Course}
A normal and successful postoperative course following circumcision is characterized by a predictable healing process with minimal discomfort and the resolution of the underlying indications. Patients can expect mild to moderate pain for the first few days, which is typically well-controlled with oral analgesics. Swelling and bruising of the penis are common and usually subside within 1-2 weeks. For neonates, parents will observe a yellowish exudate on the glans, which is normal fibrin formation, and the wound edges will gradually approximate. For adults, absorbable sutures will dissolve over 2-4 weeks, and the wound will epithelialize, forming a neat scar just proximal to the coronal sulcus.

Urination should remain normal and unobstructed throughout the recovery period, though some temporary discomfort may be experienced. Patients should gradually resume normal activities, avoiding strenuous exercise and sexual activity for the recommended period (typically 4-6 weeks for adults) to prevent complications. The primary goal is for the glans to be permanently exposed, with a satisfactory cosmetic appearance and no functional impairment. For therapeutic circumcisions, the expected course includes the complete resolution of pre-existing symptoms such as pain from phimosis, recurrent infections, or difficulty with hygiene, leading to improved quality of life.

\section*{Postoperative Care \& Follow-up}
Postoperative care for circumcision focuses on wound management, pain control, and monitoring for complications, with follow-up tailored to the patient's age.

For **neonates**:
\begin{itemize}
  \item **Parental Observation:** Parents are instructed to monitor the wound for signs of infection (excessive redness, swelling, pus, foul-smelling discharge), persistent bleeding, or difficulty urinating.
  \item **Dressing Changes:** Petroleum jelly-impregnated gauze should be applied with each diaper change for 3-5 days to prevent the glans from sticking to the diaper.
  \item **Routine Pediatric Check-up:** The infant's pediatrician will typically assess the healing at the first routine well-baby visit, usually within 1-2 weeks post-procedure.
  \item **Warning Signs:** Parents should contact the doctor immediately for fever, excessive crying, persistent bleeding (more than a few drops), foul-smelling discharge, or inability to urinate for several hours.
\end{itemize}

For **older children and adults**:
\begin{itemize}
  \item **Wound Care:** Patients are advised to keep the wound clean and dry. Daily gentle washing with soap and water is usually sufficient after the initial dressing is removed (typically 24-48 hours).
  \item **Pain Management:** Oral analgesics should be taken as prescribed.
  \item **Activity Restrictions:** Avoid strenuous exercise, heavy lifting, and sexual activity for 4-6 weeks to prevent wound dehiscence or bleeding.
  \item **Post-operative Clinic Visit:** A follow-up appointment is usually scheduled 1-2 weeks after surgery to assess wound healing, remove any non-absorbable sutures (if used), and address any concerns.
  \item **Warning Signs:** Patients are advised to seek immediate medical attention for signs of infection (increasing pain, redness, swelling, pus, fever), significant bleeding, severe pain unresponsive to medication, or difficulty urinating.
\end{itemize}
Long-term follow-up is generally not required unless complications develop or the original medical condition (e.g., BXO, penile cancer) necessitates ongoing monitoring.

\section*{Epidemiology}
Circumcision is one of the most common surgical procedures globally, with an estimated worldwide prevalence of approximately 30-33\% among males. However, its incidence varies dramatically by geographic region, culture, and religion. The highest rates are found in the Middle East, North Africa, and parts of West Africa, where it is almost universally practiced due to religious (Islam, Judaism) and cultural traditions. In the United States, the prevalence of male circumcision has historically been high, fluctuating between 58\% and 81\% in recent decades, primarily driven by non-religious factors and parental choice, though rates have shown a slight decline in recent years. In contrast, rates are very low in many European countries, parts of Asia (e.g., China, Japan), and Latin America. Medically indicated circumcisions, such as for pathological phimosis or balanitis, represent a smaller but significant proportion of procedures, particularly in adult males, and are performed globally irrespective of cultural background. Morbidity rates are generally low, estimated at 0.2\% to 3\% for neonates and slightly higher for adults, with bleeding and infection being the most common complications. Mortality is exceedingly rare, estimated at less than 1 in 500,000 procedures.

\section*{Outcomes \& Prognosis}
The outcomes of circumcision are generally excellent, with high success rates in achieving the intended goals, whether performed electively or therapeutically. For elective neonatal circumcisions, the prognosis is overwhelmingly positive, with minimal long-term complications and high parental satisfaction regarding hygiene and cosmetic appearance. When performed for medical indications such as pathological phimosis, paraphimosis, or recurrent balanoposthitis, the procedure typically provides definitive resolution of symptoms, with success rates often exceeding 95\%. Patients experience improved hygiene, reduced incidence of local infections, and relief from pain or discomfort.

Long-term studies suggest that circumcision may offer prophylactic benefits, including a reduced risk of urinary tract infections in infancy, particularly in boys with urinary tract anomalies. Furthermore, several randomized controlled trials (e.g., those conducted in sub-Saharan Africa) have demonstrated that male circumcision significantly reduces the risk of acquiring certain sexually transmitted infections, notably HIV, human papillomavirus (HPV), and herpes simplex virus (HSV), in adulthood, particularly in high-prevalence regions. The risk of penile cancer, while rare, is also significantly lower in circumcised males. Prognostic factors for a successful outcome include meticulous surgical technique, appropriate patient selection (avoiding contraindications), and adherence to postoperative care instructions. Recurrence of foreskin issues, such as skin bridges or meatal stenosis, can occur but are generally amenable to minor revision procedures. Overall, the long-term functional and cosmetic outcomes are favorable for the vast majority of individuals undergoing the procedure.

\section*{References}
\begin{enumerate}
  \item MedlinePlus. Circumcision. U.S. National Library of Medicine; 2024. Available from: \url{https://medlineplus.gov/circumcision.html}
  \item Wein AJ, Kavoussi LR, Partin AW, Peters CA, editors. Campbell-Walsh-Wein Urology. 12th ed. Philadelphia, PA: Elsevier; 2021.
  \item American Academy of Pediatrics Task Force on Circumcision. Male Circumcision. Pediatrics. 2012;130(3):e756-e785.
  \item Sabiston DC, Townsend CM, Beauchamp RD, Evers BM, Mattox KL, editors. Sabiston Textbook of Surgery: The Biological Basis of Modern Surgical Practice. 21st ed. Philadelphia, PA: Elsevier; 2022.
\end{enumerate}

\clearpage